\section{Fouille de Données}

\begin{frame}
\frametitle{Introduction à la Fouille de Données}
\framesubtitle{Data Mining}

\textbf{Définition}

\vspace{0.3cm}

La fouille de données est le processus d'\textbf{extraction de connaissances} implicites, utiles et compréhensibles à partir de grandes quantités de données.

\vspace{0.3cm}

\textbf{Elle combine :}
\begin{itemize}
    \item Statistiques
    \item Intelligence artificielle (apprentissage automatique)
    \item Visualisation
    \item Algorithmes de recherche de motifs
\end{itemize}

\end{frame}

\begin{frame}{Historique du Data Mining}

\begin{itemize}
    \item<+-> \textbf{Années 1960-1970}
        \begin{itemize}
            \item Développement des bases de données relationnelles (SQL)
            \item Premières techniques d'analyse statistique
        \end{itemize}

    \vspace{0.2cm}
    \item<+-> \textbf{Années 1980-1990}
        \begin{itemize}
            \item Boom avec l'augmentation des volumes de données
            \item Augmentation de la puissance de calcul
            \item Apparition du terme "data mining" dans les années 1990
            \item Popularisé par le projet KDD (Knowledge Discovery in Databases)
        \end{itemize}

    \vspace{0.2cm}
    \item<+-> \textbf{Années 2000 à aujourd'hui}
        \begin{itemize}
            \item Intégration avec l'IA et le big data
            \item Démocratisation : outils accessibles (Python, R, Tableau)
        \end{itemize}
\end{itemize}

\end{frame}

\begin{frame}{Applications : Finance}

\textbf{Problème :} Détection de fraudes dans les transactions bancaires

\vspace{0.3cm}

\textbf{Solution :}
\begin{itemize}
    \item<+-> Clustering pour identifier les comportements inhabituels
        \begin{itemize}
            \item Exemple : dépenses soudaines anormales
        \end{itemize}
    \item<+-> Modèles de classification pour prédire si une transaction est frauduleuse
\end{itemize}

\vspace{0.3cm}

\onslide<2>
\textbf{Algorithmes utilisés :}
\begin{itemize}
    \item Arbres de décision
    \item Réseaux de neurones
\end{itemize}

\end{frame}

\begin{frame}{Applications : Santé}

\textbf{Problème :} Prédiction des risques de maladies chroniques

\vspace{0.3cm}

\textbf{Solution :}
\begin{itemize}
    \item<+-> Analyse de dossiers médicaux
        \begin{itemize}
            \item Identifier des corrélations entre symptômes et diagnostics
        \end{itemize}
    \item<+-> Visualisation des résultats pour aider les médecins
        \begin{itemize}
            \item Exemple : heatmaps des zones à risque dans une population
        \end{itemize}
\end{itemize}

\end{frame}

\begin{frame}{Applications : Marketing}

\textbf{Problème :} Segmentation des clients pour personnaliser les offres

\vspace{0.3cm}

\textbf{Solution :}
\begin{itemize}
    \item<+-> \textbf{Clustering (k-means)}
        \begin{itemize}
            \item Regrouper les clients selon leurs comportements d'achat
        \end{itemize}

    \vspace{0.2cm}
    \item<+-> \textbf{Règles d'association}
        \begin{itemize}
            \item Identifier des produits souvent achetés ensemble
        \end{itemize}
\end{itemize}

\end{frame}

\begin{frame}{Applications : Retail et Industrie}

\begin{columns}
\begin{column}{0.48\textwidth}
    \colormark{blue}\textbf{Retail :}
    \begin{itemize}
        \item \textbf{Problème :} Optimisation des stocks et prévision des ventes
        \item \textbf{Solution :}
        \begin{itemize}
            \item Séries temporelles pour prédire la demande
            \item Analyse de paniers d'achat
        \end{itemize}
    \end{itemize}
\end{column}

\begin{column}{0.48\textwidth}
    \onslide<2>
    \colormark{blue}\textbf{Industrie :}
    \begin{itemize}
        \item \textbf{Problème :} Maintenance prédictive des machines
        \item \textbf{Solution :}
        \begin{itemize}
            \item Analyse des données de capteurs
            \item Réduction de dimensionnalité (PCA)
        \end{itemize}
    \end{itemize}
\end{column}
\end{columns}

\end{frame}

\begin{frame}{Méthodologie : CRISP-DM}
\framesubtitle{Cross-Industry Standard Process for Data Mining}

\textbf{Une approche structurée en 6 étapes :}

\vspace{0.3cm}

\begin{enumerate}
    \item Compréhension du problème métier
    \item Compréhension des données
    \item Préparation des données
    \item Modélisation
    \item Évaluation
    \item Déploiement et maintenance
\end{enumerate}

\end{frame}

\begin{frame}{CRISP-DM : 1. Compréhension du problème}

\textbf{Objectif :} Définir clairement le problème à résoudre

\vspace{0.3cm}

\begin{itemize}
    \item Exemple : "Comment réduire le taux d'attrition des clients ?"
    \item Impliquer les parties prenantes
    \item Définir un objectif clair et actionnable
\end{itemize}

\vspace{0.5cm}

\onslide<2>
\textbf{Importance de la formulation de la question :}

\vspace{0.2cm}

\begin{itemize}
    \item \textcolor{red}{Mauvaise question :} "Quelles sont les tendances de nos ventes ?"
        \begin{itemize}
            \item Trop vague, difficile à actionner
        \end{itemize}
    \item \textcolor{green}{Bonne question :} "Quels produits ont un taux de retour élevé et quels facteurs (taille, couleur, saison) y contribuent ?"
        \begin{itemize}
            \item Permet d'identifier des leviers d'amélioration
        \end{itemize}
\end{itemize}

\end{frame}

\begin{frame}{CRISP-DM : 2. Compréhension des données}

\textbf{Exploration (EDA) :}

\vspace{0.3cm}

\begin{itemize}
    \item<+-> Analyser la distribution des données
    \item<+-> Identifier les valeurs manquantes ou aberrantes
    \item<+-> Visualisation : graphiques
        \begin{itemize}
            \item Histogrammes
            \item Boxplots
            \item Heatmaps
        \end{itemize}
\end{itemize}

\end{frame}

\begin{frame}{CRISP-DM : 3. Préparation des données}

\textbf{Nettoyage :}

\vspace{0.3cm}

\begin{itemize}
    \item<+-> Gestion des valeurs manquantes
        \begin{itemize}
            \item Imputation
            \item Suppression
        \end{itemize}

    \vspace{0.2cm}
    \item<+-> Normalisation
        \begin{itemize}
            \item Mise à l'échelle pour algorithmes (ex : k-means)
        \end{itemize}
\end{itemize}

\vspace{0.3cm}

\onslide<2>
\textbf{Transformation :}
\begin{itemize}
    \item Création de nouvelles variables
    \item Exemples : ratio, catégories
\end{itemize}

\end{frame}

\begin{frame}{CRISP-DM : 4. Modélisation}

\textbf{Choix du modèle en fonction du problème :}

\vspace{0.3cm}

\begin{itemize}
    \item<+-> \textbf{Classification} : prédire une catégorie
        \begin{itemize}
            \item Exemple : spam/non-spam
            \item Algorithmes : Arbres de décision, SVM, Réseaux de neurones
        \end{itemize}

    \vspace{0.2cm}
    \item<+-> \textbf{Régression} : prédire une valeur continue
        \begin{itemize}
            \item Exemple : prix d'une maison
        \end{itemize}

    \vspace{0.2cm}
    \item<+-> \textbf{Clustering} : regrouper des données sans étiquettes
        \begin{itemize}
            \item Exemple : segmentation client
        \end{itemize}

    \vspace{0.2cm}
    \item<+-> \textbf{Règles d'association} : trouver des motifs
        \begin{itemize}
            \item Exemple : "30\% des clients qui achètent X achètent aussi Y"
        \end{itemize}
\end{itemize}

\end{frame}

\begin{frame}{CRISP-DM : 5. Évaluation}

\textbf{Métriques adaptées :}

\vspace{0.3cm}

\begin{itemize}
    \item<+-> \textbf{Classification :}
        \begin{itemize}
            \item Précision
            \item Rappel
            \item F1-score
        \end{itemize}

    \vspace{0.2cm}
    \item<+-> \textbf{Régression :}
        \begin{itemize}
            \item RMSE (Root Mean Squared Error)
        \end{itemize}
\end{itemize}

\vspace{0.3cm}

\onslide<2>
\textbf{Validation croisée :}
\begin{itemize}
    \item Pour éviter le surfitement (overfitting)
\end{itemize}

\end{frame}

\begin{frame}{CRISP-DM : 6. Déploiement}

\textbf{Déploiement et maintenance :}

\vspace{0.3cm}

\begin{itemize}
    \item<+-> Intégration du modèle dans un système opérationnel
    \item<+-> Surveillance continue
        \begin{itemize}
            \item Détecter la dérive des données
            \item Détecter la perte de performance
        \end{itemize}
\end{itemize}

\end{frame}

\begin{frame}{Limites : Problèmes courants}

\begin{enumerate}
    \item<+-> \textbf{Mauvaise formulation de la question}
        \begin{itemize}
            \item Risque : résultats inutiles ou non actionnables
            \item Solution : Co-construction avec les métiers
        \end{itemize}

    \vspace{0.2cm}
    \item<+-> \textbf{Données insuffisantes ou de mauvaise qualité}
        \begin{itemize}
            \item Problème du "rare event" : cas rares nécessitent beaucoup de données
            \item Biais dans les données
            \item Exemple : modèle de recrutement biaisé
        \end{itemize}

    \vspace{0.2cm}
    \item<+-> \textbf{Surfitement (overfitting)}
        \begin{itemize}
            \item Le modèle mémorise mais ne généralise pas
            \item Solution : validation croisée, régularisation
        \end{itemize}
\end{enumerate}

\end{frame}

\begin{frame}{Limites : Interprétabilité}

\textbf{Problème :}

\vspace{0.3cm}

\begin{itemize}
    \item Certains algorithmes sont des "boîtes noires"
    \item Exemple : réseaux de neurones profonds
\end{itemize}

\vspace{0.5cm}

\onslide<2>
\textbf{Solution :}
\begin{itemize}
    \item Privilégier des modèles explicables quand c'est possible
    \item Arbres de décision
    \item Régression linéaire
\end{itemize}

\end{frame}

\begin{frame}{Bonnes pratiques}

\small
\begin{tabular}{|p{3.5cm}|p{3cm}|p{3cm}|}
\hline
\textbf{Bonne pratique} & \textbf{Description} & \textbf{Exemple} \\
\hline
Question claire et actionnable & Impliquer les parties prenantes & "Réduire coûts logistiques de 10\% sans impacter satisfaction" \\
\hline
Qualité des données & Nettoyer et valider avant analyse & Supprimer doublons, imputer valeurs manquantes \\
\hline
Validation robuste & Cross-validation pour généralisation & Diviser données en 5 folds \\
\hline
Documentation & Documenter chaque étape & Notebooks Jupyter commentés \\
\hline
Éthique et transparence & Auditer pour éviter biais & Vérifier non-discrimination \\
\hline
\end{tabular}

\end{frame}

\begin{frame}{Collecte de données : Sources}

\textbf{1. Données structurées}
\begin{itemize}
    \item Bases de données relationnelles (SQL)
    \item Fichiers CSV/Excel
    \item Exemple : bases clients, transactions
\end{itemize}

\vspace{0.3cm}

\onslide<2->
\textbf{2. Données non structurées}
\begin{itemize}
    \item Textes (emails, avis clients)
    \item Images, vidéos
    \item Exemple : analyse de sentiments (text mining)
\end{itemize}

\vspace{0.3cm}

\onslide<3>
\textbf{3. Données externes}
\begin{itemize}
    \item Open data (INSEE, Eurostat)
    \item APIs (météo, trafic)
    \item Web scraping (respecter RGPD)
\end{itemize}

\end{frame}

\begin{frame}{Collecte de données : Défis}

\textbf{Accès et légalité :}
\begin{itemize}
    \item Respecter les réglementations (RGPD)
    \item Obtenir les droits d'utilisation
\end{itemize}

\vspace{0.3cm}

\onslide<2->
\textbf{Nettoyage et préparation :}
\begin{itemize}
    \item Données bruitées, incohérentes, manquantes
    \item Techniques : imputation, suppression des doublons, normalisation
\end{itemize}

\vspace{0.3cm}

\onslide<3>
\textbf{Intégration :}
\begin{itemize}
    \item Fusionner des données de sources différentes
    \item Exemple : données internes + données marché
\end{itemize}

\end{frame}

\begin{frame}{Algorithmes spécifiques au data mining}

\begin{enumerate}
    \item<+-> \textbf{Règles d'association}
        \begin{itemize}
            \item Découvrir des associations entre variables
            \item Exemple : catégorie socio-professionnelle vs certains produits
            \item Utilisé en marketing pour recommandations
        \end{itemize}

    \vspace{0.3cm}
    \item<+-> \textbf{Détection d'anomalies}
        \begin{itemize}
            \item Identifier des observations rares ou suspectes
            \item Exemple : fraudes bancaires
            \item Techniques : isolation forest, auto-encoders
        \end{itemize}
\end{enumerate}

\end{frame}

\begin{frame}{Cas pratique : Optimisation supermarché}
\framesubtitle{Réduire le gaspillage et optimiser les promotions}

\textbf{Scénario :}

\vspace{0.2cm}

Vous travaillez pour un supermarché qui souhaite :
\begin{itemize}
    \item Optimiser ses promotions
    \item Réduire le gaspillage alimentaire
\end{itemize}

\vspace{0.3cm}

\textbf{Comment aborderiez-vous ce problème ?}

\end{frame}

\begin{frame}{Cas pratique : 1. Formulation de la question}

\textbf{Questions à poser :}

\vspace{0.3cm}

\begin{itemize}
    \item "Quels produits ont un taux de retour élevé (gaspillage) ?"
    \item "Quels facteurs y contribuent ?"
        \begin{itemize}
            \item Prix
            \item Saison
            \item Emplacement en rayon
        \end{itemize}
    \item "Quelles promotions croisées maximisent les ventes ?"
        \begin{itemize}
            \item Exemple : fromage + vin
            \item Sans impacter la marge
        \end{itemize}
\end{itemize}

\end{frame}

\begin{frame}{Cas pratique : 2-3. Collecte et Préparation}

\textbf{Collecte des données :}
\begin{itemize}
    \item Données internes : ventes, stocks, promotions passées
    \item Données externes : météo, événements locaux
\end{itemize}

\vspace{0.3cm}

\onslide<2>
\textbf{Préparation :}
\begin{itemize}
    \item Nettoyage : supprimer transactions incomplètes
    \item Imputer valeurs manquantes
    \item Créer nouvelles variables :
        \begin{itemize}
            \item Ratio gaspillage = (stock final - ventes) / stock initial
        \end{itemize}
\end{itemize}

\end{frame}

\begin{frame}{Cas pratique : 4. Modélisation}

\textbf{Pour le gaspillage :}
\begin{itemize}
    \item<+-> Régression pour prédire le taux de gaspillage
        \begin{itemize}
            \item En fonction du prix et de la saison
        \end{itemize}
    \item<+-> Clustering pour segmenter les produits
        \begin{itemize}
            \item Selon sensibilité au gaspillage
        \end{itemize}
\end{itemize}

\vspace{0.3cm}

\onslide<3>
\textbf{Pour les promotions croisées :}
\begin{itemize}
    \item Règles d'association
        \begin{itemize}
            \item Trouver produits souvent achetés ensemble
        \end{itemize}
\end{itemize}

\end{frame}

\begin{frame}{Cas pratique : 5-6. Évaluation et Déploiement}

\textbf{Évaluation :}
\begin{itemize}
    \item Pour la régression : RMSE pour mesurer l'erreur
    \item Pour les règles d'association : support, confidence, lift
\end{itemize}

\vspace{0.3cm}

\onslide<2>
\textbf{Déploiement :}
\begin{itemize}
    \item Recommander promotions ciblées
        \begin{itemize}
            \item "Achetez le fromage et obtenez 20\% sur le vin"
        \end{itemize}
    \item Mettre en place un tableau de bord
        \begin{itemize}
            \item Suivre l'impact des promotions
        \end{itemize}
\end{itemize}

\end{frame}

\begin{frame}{Outils et ressources}

\textbf{Pour la visualisation :}
\begin{itemize}
    \item Tableau, Power BI (drag-and-drop)
    \item Idéal pour les non-techniciens
\end{itemize}

\vspace{0.3cm}

\onslide<2->
\textbf{Pour le nettoyage et l'analyse :}
\begin{itemize}
    \item Excel (formules avancées, Power Query)
    \item Python : Pandas, Matplotlib/Seaborn, Scikit-learn
    \item R : analyses statistiques avancées
\end{itemize}

\end{frame}

\begin{frame}{Conclusion : Points clés}

\textbf{Bonnes pratiques à retenir :}

\vspace{0.3cm}

\begin{enumerate}
    \item \textbf{Poser la bonne question}

    Impliquer les métiers, objectif clair et actionnable

    \item \textbf{Travailler avec des données de qualité}
    
    Nettoyage rigoureux, gestion des valeurs manquantes

    \item \textbf{Évaluer correctement les modèles}
    
    Validation croisée, métriques adaptées

    \item \textbf{Documenter et expliquer}
    
    Modèle compréhensible par tous

    \item \textbf{Penser à l'éthique}
    
    Auditer les données pour éviter les biais

\end{enumerate}

\end{frame}

\begin{frame}{Message final}

\textbf{Le data mining n'est pas une magie}

\vspace{0.3cm}

C'est un processus structuré qui repose sur :

\vspace{0.3cm}

\begin{itemize}
    \item Une \textcolor{blue}{bonne formulation du problème}
    \item Des \textcolor{blue}{données propres et représentatives}
    \item Des \textcolor{blue}{méthodes adaptées}
    \item Et une \textcolor{blue}{évaluation rigoureuse}
\end{itemize}

\vspace{0.5cm}

En suivant ces bonnes pratiques, même sans expertise technique avancée, il est possible de tirer des insights puissants pour innover et prendre des décisions éclairées.

\end{frame}
